% !TEX root = ../main.tex
%% Corrigés — Chapitre 16 — Probabilités conditionnelles
%% Tous les résultats ont été vérifiés par programme (calcul exact en fractions et énumération des issues).

\corrigesChapitre{16}{Probabilités conditionnelles}

\begin{deuxcolonnes}

\begin{corrige}{1}{Arbres pondérés}
\methode{produit le long d'un chemin, somme des chemins}
\begin{questions}
  \item $\frac35\times\frac24=\res{\frac{3}{10}}$
  \item Chemins PF et FP : $2\times\frac14=\res{\frac12}$
  \item Chemins BN et NB : $2\times0{,}4\times0{,}6=\res{\frac{12}{25}}$
  \item $\left(\frac25\right)^{2}+\left(\frac25\right)^{2}+\left(\frac15\right)^{2}=\res{\frac{9}{25}}$
  \item Par le contraire (trois bleues) :\\ $1-\frac{6}{10}\times\frac59\times\frac48=1-\frac16=\res{\frac56}$
\end{questions}
\end{corrige}

\begin{corrige}{2}{Calculs de base}
\begin{questions}
  \item $P_B(A)=\frac{0{,}2}{0{,}4}=\res{0{,}5}$ et\\ $P_A(B)=\frac{0{,}2}{0{,}6}=\res{\frac13}$
  \item Somme supérieure à $9$ : $6$ issues équiprobables, dont $3$ commencent par $6$ : $\res{\frac12}$
  \item Il reste $3$ rouges sur $9$ boules : $\res{\frac13}$
  \item $P(A\cap B)=P(B)\times P_B(A)=\res{0{,}21}$
  \item $3$ cœurs parmi les $12$ figures : $\res{\frac14}$
\end{questions}
\end{corrige}

\begin{corrige}{3}{Indépendance}
\begin{questions}
  \item $0{,}4\times0{,}5=0{,}2=P(A\cap B)$ :\\ $\res{\text{indépendants}}$
  \item $\frac{1}{36}=\frac16\times\frac16$ : $\res{\text{indépendants}}$
  \item $0{,}3+0{,}6-0{,}3\times0{,}6=\res{0{,}72}$
  \item Oui : avec remise, le second tirage se fait dans les mêmes conditions.
  \item $P(A)\,P\big(\overline B\big)=0{,}5\times0{,}6=\res{0{,}3}$
  \item $A\cap B=\varnothing$ : $\res{\text{incompatibles}}$. Mais $P(A)P(B)=\frac1{12}\neq0$ : $\res{\text{pas indépendants}}$.
\end{questions}
\end{corrige}

\begin{corrige}{4}{Arbres pondérés et applications}
\begin{questions}
  \item $\frac58\times\frac37=\res{\frac{15}{56}}$
  \item Somme $10$ : $(4\pv6)$, $(5\pv5)$, $(6\pv4)$ : $\res{\frac13}$
  \item $\frac{18}{30}\times\frac{10}{18}+\frac{12}{30}\times\frac{8}{12}$\\ $=\frac{18}{30}=\res{\frac35}$
  \item $0{,}05\times0{,}98+0{,}95\times0{,}05$\\ $=\res{0{,}096\,5}$
  \item $3$ as parmi $51$ cartes : $\res{\frac{1}{17}}$
\end{questions}
\end{corrige}

\begin{corrige}{5}{Indépendance et applications}
\begin{questions}
  \item $P=\frac{9}{36}=\frac14=\frac12\times\frac12$ : $\res{\text{oui}}$
  \item $0{,}6^{2}=\res{\frac{9}{25}}$
  \item PPF, PFP, FPP : $\res{\frac38}$
  \item $0{,}4+0{,}6-0{,}24=\res{0{,}76}$
  \item $0{,}9\times0{,}8=\res{0{,}72}$
\end{questions}
\end{corrige}

\begin{corrige}{6}{Formule des probabilités totales}
\begin{questions}
  \item $0{,}008+0{,}010\,5+0{,}012\,5=\res{0{,}031}$
  \item $0{,}6\times0{,}7+0{,}4\times0{,}5=\res{0{,}62}$
  \item $0{,}02\times0{,}95+0{,}98\times0{,}10$\\ $=0{,}019+0{,}098=\res{0{,}117}$
  \item $\frac12\times\frac35+\frac12\times\frac15=\res{\frac25}$
  \item $0{,}63+0{,}09=\res{0{,}72}$
\end{questions}
\end{corrige}

\begin{corrige}{7}{Inverser un conditionnement}
\methode{$P_B(A)=\frac{P(A\cap B)}{P(B)}$}
\begin{questions}
  \item $\frac{0{,}012\,5}{0{,}031}\approx\res{0{,}403}$
  \item $\frac{0{,}019}{0{,}117}\approx\res{0{,}162}$ : un test positif ne signifie pas que l'on est malade ; la maladie est rare
    et les « faux positifs » sont nombreux.
  \item $\frac{0{,}3}{0{,}4}=\res{\frac34}$
  \item $\frac{0{,}63}{0{,}72}=\res{0{,}875}$
\end{questions}
\end{corrige}

\end{deuxcolonnes}

%% ---------------------------------------------------------------------
\begin{corrige}{8}{Tirages successifs sans remise}
\noindent
\begin{minipage}[t]{0.5\linewidth}\vspace{0pt}\raggedright
\begin{questions}
  \item $\frac58\times\frac47\times\frac36=\res{\frac{5}{28}}$.
  \item Trois chemins (RRB, RBR, BRR), chacun de probabilité $\frac{5\times4\times3}{8\times7\times6}=\frac{5}{28}$ : $\res{\frac{15}{28}}$.
  \item Contraire de « trois rouges » : $1-\frac{5}{28}=\res{\frac{23}{28}}$.
  \item Il reste $4$ rouges et $3$ bleues : $\frac37\times\frac26=\res{\frac17}$.
  \item Chemins RRR, RBR, BRR et BBR : $\frac{60+60+60+30}{336}=\frac{210}{336}=\res{\frac58}$, comme pour la première boule.
\end{questions}
\end{minipage}\hfill
\begin{minipage}[t]{0.46\linewidth}\vspace{0pt}
\centering
\begin{tikzpicture}[otArbre, level distance=17mm, level 1/.style={sibling distance=22mm},
  level 2/.style={sibling distance=11mm}, level 3/.style={sibling distance=5.5mm}]
  \node {}
    child {node {R}
      child {node {R}
        child {node {R} edge from parent node[otPoids] {$\frac36$}}
        child {node {B} edge from parent node[otPoids] {$\frac36$}}
        edge from parent node[otPoids] {$\frac47$}}
      child {node {B}
        child {node {R} edge from parent node[otPoids] {$\frac46$}}
        child {node {B} edge from parent node[otPoids] {$\frac26$}}
        edge from parent node[otPoids] {$\frac37$}}
      edge from parent node[otPoids] {$\frac58$}}
    child {node {B}
      child {node {R}
        child {node {R} edge from parent node[otPoids] {$\frac46$}}
        child {node {B} edge from parent node[otPoids] {$\frac26$}}
        edge from parent node[otPoids] {$\frac57$}}
      child {node {B}
        child {node {R} edge from parent node[otPoids] {$\frac56$}}
        child {node {B} edge from parent node[otPoids] {$\frac16$}}
        edge from parent node[otPoids] {$\frac27$}}
      edge from parent node[otPoids] {$\frac38$}};
\end{tikzpicture}
\end{minipage}
\end{corrige}

%% ---------------------------------------------------------------------
\begin{corrige}{9}{Deux tirages avec remise}
\begin{questions}
  \item $4\times4=\res{16}$ résultats équiprobables, dont $4$ doubles : $\res{\frac14}$.
  \item $(1\pv4)$, $(2\pv3)$, $(3\pv2)$, $(4\pv1)$ : $\res{\frac14}$.
  \item Le produit est impair si les deux numéros sont impairs ($4$ issues) : $1-\frac{4}{16}=\res{\frac34}$.
  \item La somme est paire pour $8$ issues (deux pairs ou deux impairs), dont $4$ avec deux numéros pairs : $\res{\frac12}$.
\end{questions}
\end{corrige}

%% ---------------------------------------------------------------------
\begin{corrige}{10}{Applications complexes}
\begin{questions}
  \item $P(D)=0{,}5\times0{,}01+0{,}3\times0{,}02+0{,}2\times0{,}04=\res{0{,}019}$ et $P_D(C)=\frac{0{,}008}{0{,}019}\approx\res{0{,}421}$.
  \item Tirage simultané équivalent : $\frac{\binom42\binom{48}1}{\binom{52}3}=\frac{288}{22\,100}\approx\res{0{,}013}$.
  \item $P(T)=0{,}03\times0{,}96+0{,}97\times0{,}06=0{,}028\,8+0{,}058\,2=\res{0{,}087}$ et $P_T(M)=\frac{0{,}028\,8}{0{,}087}\approx\res{0{,}331}$.
  \item Avec remise, les tirages sont indépendants : $1-\left(\frac58\right)^{3}=\res{\frac{387}{512}\approx0{,}756}$, contre
    $\frac{23}{28}\approx0{,}821$ sans remise : sans remise, chaque rouge tirée rend une bleue plus probable ensuite.
  \item $1-\left(\frac56\right)^{4}=\res{\frac{671}{1\,296}\approx0{,}518}$.
\end{questions}
\end{corrige}

%% ---------------------------------------------------------------------
\begin{corrige}{11}{Devoir de synthèse}
\begin{questions}
  \item $P_B(A)=\frac{0{,}15}{0{,}4}=\res{0{,}375}$ et $P_A(B)=\frac{0{,}15}{0{,}5}=\res{0{,}3}$.
  \item $\res{\frac{1}{13}}$ (un as parmi les $13$ cœurs).
  \item Chemins RB et BR : $2\times\frac{4}{10}\times\frac69=\res{\frac{8}{15}}$.
  \item $\res{\frac38}$.
  \item $P(\text{les deux})=\frac{3\times3}{36}=\frac14=\frac12\times\frac12$ : $\res{\text{indépendants}}$.
  \item $0{,}6+0{,}5-0{,}3=\res{0{,}8}$.
  \item $P(D)=0{,}006+0{,}012=\res{0{,}018}$ et $P_D(M_2)=\frac{0{,}012}{0{,}018}=\res{\frac23}$.
  \item $0{,}55\times0{,}75+0{,}45\times0{,}6=\res{0{,}682\,5}$.
  \item $P_T(M)=\frac{0{,}04\times0{,}97}{0{,}04\times0{,}97+0{,}96\times0{,}08}=\frac{0{,}038\,8}{0{,}115\,6}\approx\res{0{,}336}$.
\end{questions}
\end{corrige}

%% ---------------------------------------------------------------------
\begin{corrige}{12}{Le paradoxe des tests de dépistage}
\partieExo{A}{Une maladie rare : $p=0{,}01$}
\noindent
\begin{minipage}[t]{0.5\linewidth}\vspace{0pt}\raggedright
\begin{questions}
  \item Sur $10\,000$ personnes, $100$ sont malades : $99$ ont un test positif. Parmi les $9\,900$ personnes saines, 2\,\% ont un test
    positif, soit $198$ (tableau ci-contre).
  \item $P(T)=\frac{297}{10\,000}=\res{0{,}029\,7}$ et $P_T(M)=\frac{99}{297}=\res{\frac13}$ : deux tests positifs sur trois sont des
    \emph{faux positifs}, car les personnes saines sont bien plus nombreuses.
\end{questions}
\end{minipage}\hfill
\begin{minipage}[t]{0.46\linewidth}\vspace{1mm}
\centering
{\renewcommand{\arraystretch}{1.5}\arrayrulecolor{otRule}%
\begin{tabular}{@{}l>{\centering\arraybackslash}p{13mm}>{\centering\arraybackslash}p{13mm}>{\centering\arraybackslash}p{13mm}@{}}
\toprule
 & $T$ & $\overline T$ & Total \\
\midrule
$M$ & $99$ & $1$ & $100$ \\
$\overline M$ & $198$ & $9\,702$ & $9\,900$ \\
\midrule
Total & $297$ & $9\,703$ & $10\,000$ \\
\bottomrule
\end{tabular}}
\end{minipage}
\begin{questions}[start=3]
  \item $P(M\cap T_1\cap T_2)=0{,}01\times0{,}99^{2}=0{,}009\,801$ et $P\big(\overline M\cap T_1\cap T_2\big)=0{,}99\times0{,}02^{2}=0{,}000\,396$. D'où la
    probabilité $\frac{0{,}009\,801}{0{,}009\,801+0{,}000\,396}\approx\res{0{,}961}$ : le second test est très efficace.
\end{questions}
\partieExo{B}{Influence de la proportion de malades}
\noindent
\begin{minipage}[t]{0.54\linewidth}\vspace{0pt}\raggedright
\begin{questions}[start=4]
  \item Probabilités totales : $P(T)=0{,}99\,p+0{,}02\,(1-p)=\res{0{,}97\,p+0{,}02}$. Puis $P_T(M)=\frac{0{,}99\,p}{0{,}97\,p+0{,}02}=\frac{99\,p}{97\,p+2}$
    (on multiplie par $100$).
  \item $f'(p)=\frac{99(97p+2)-99p\times97}{(97p+2)^{2}}=\frac{198}{(97p+2)^{2}}>0$ : $\res{f\text{ est croissante}}$ sur $\intOF{0}{1}$, avec $f(1)=1$.
  \item $f(p)\geq0{,}9\iff99p\geq87{,}3\,p+1{,}8\iff11{,}7\,p\geq1{,}8\iff p\geq\frac{2}{13}$ : il faut $\res{p\geq\frac{2}{13}\approx0{,}154}$,
    soit au moins $15{,}4$\,\% de malades.
\end{questions}
\end{minipage}\hfill
\begin{minipage}[t]{0.42\linewidth}\vspace{1mm}
\centering
\begin{tikzpicture}
\begin{axis}[otSchema, width=66mm, height=46mm, xmin=0, xmax=1.05, ymin=0, ymax=1.08,
  xtick={0.154,0.5,1}, xticklabels={$\frac{2}{13}$,$0{,}5$,$1$}, ytick={0.9,1}, yticklabels={$0{,}9$,$1$},
  xlabel={$p$}, xlabel style={at={(1,0)},anchor=west}, grid=major, grid style={otRule!70,line width=0.3pt}]
  \addplot[otNavy,line width=1.2pt,domain=0.001:1,samples=200] {99*x/(97*x+2)};
  \draw[otGreen,dashed,line width=0.6pt] (axis cs:0,0.9) -- (axis cs:1,0.9);
  \fill[otGreen] (axis cs:0.1538,0.9) circle (1.5pt);
\end{axis}
\end{tikzpicture}\\[0.5mm]
\legende{$P_T(M)$ en fonction de la proportion $p$ de malades.}
\end{minipage}
\end{corrige}

%% ---------------------------------------------------------------------
\begin{corrige}{13}{Une suite de probabilités}
\noindent
\begin{minipage}[t]{0.56\linewidth}\vspace{0pt}\raggedright
\begin{questions}
  \item $p_2=\res{0{,}7}$ ; $p_3=0{,}7\times0{,}7+0{,}3\times0{,}4=\res{0{,}61}$.
  \item D'après l'arbre et la formule des probabilités totales : $p_{n+1}=0{,}7\,p_n+0{,}4\,(1-p_n)=\res{0{,}3\,p_n+0{,}4}$.
  \item $u_{n+1}=0{,}3\,p_n+0{,}4-\frac47=0{,}3\,p_n-\frac{6}{35}=0{,}3\left(p_n-\frac47\right)=0{,}3\,u_n$ : la suite est géométrique de
    raison $0{,}3$, avec $u_1=\frac37$. Donc $\res{p_n=\frac47+\frac37\times0{,}3^{\,n-1}}$.
  \item $0<0{,}3<1$, donc $0{,}3^{\,n-1}\to0$ et $\res{\lim p_n=\frac47\approx0{,}571}$ : à long terme, Léa prend le vélo environ
    $4$ jours sur $7$, quel que soit son choix du premier jour.
  \item La fonction renvoie $\res{7}$ : à partir du $7^{\text{e}}$ jour, $p_n$ est à moins de $10^{-3}$ de sa limite.
\end{questions}
\end{minipage}\hfill
\begin{minipage}[t]{0.4\linewidth}\vspace{2mm}
\centering
\begin{tikzpicture}[otArbre, level distance=20mm]
  \node {}
    child {node {$V_n$}
      child {node {$V_{n+1}$} edge from parent node[otPoids] {$0{,}7$}}
      child {node {$\overline{V_{n+1}}$} edge from parent node[otPoids] {$0{,}3$}}
      edge from parent node[otPoids] {$p_n$}}
    child {node {$\overline{V_n}$}
      child {node {$V_{n+1}$} edge from parent node[otPoids] {$0{,}4$}}
      child {node {$\overline{V_{n+1}}$} edge from parent node[otPoids] {$0{,}6$}}
      edge from parent node[otPoids] {$1-p_n$}};
\end{tikzpicture}
\end{minipage}
\end{corrige}

%% ---------------------------------------------------------------------
\begin{corrige}{14}{Fiabilité d'un système}
\begin{questions}
  \item Par indépendance : $0{,}9^{2}=\res{0{,}81}$. Puis $0{,}9^{n}<0{,}5\iff n>\frac{\ln0{,}5}{\ln0{,}9}\approx6{,}58$ (car $\ln0{,}9<0$) :
    $\res{n=7}$.
  \item L'appareil tombe en panne si les deux composants sont en panne : $1-0{,}1^{2}=\res{0{,}99}$. Puis
    $1-0{,}1^{n}\geq0{,}999\,9\iff0{,}1^{n}\leq10^{-4}\iff n\geq4$ : $\res{4\text{ composants}}$.
  \item (a) Chaque branche fonctionne avec la probabilité $0{,}81$ ; en parallèle : $1-0{,}19^{2}=\res{0{,}963\,9}$.
    (b) Chaque bloc fonctionne avec la probabilité $0{,}99$ ; en série : $0{,}99^{2}=\res{0{,}980\,1}$. Le montage (b) est le plus fiable.
  \item En notant $C_1$ et $C_2$ les événements « le composant fonctionne » :
    $P_{C_1\cup C_2}(C_1\cap C_2)=\frac{P(C_1\cap C_2)}{P(C_1\cup C_2)}=\frac{0{,}81}{0{,}99}=\res{\frac{9}{11}\approx0{,}818}$.
\end{questions}
\end{corrige}
